\title{Parameter-Efficient Orthogonal Finetuning via Factorization}

\begin{abstract}
The proliferation of large foundation models has rendered training from scratch increasingly cost-prohibitive. As a consequence, methods for efficiently adapting these models to specialized tasks have become crucial. This study investigates a systematic finetuning strategy, Orthogonal Finetuning (OFT), designed for downstream task adaptation. While OFT offers robust generalizability, its reliance on high-dimensional orthogonal matrices entails a substantial number of tunable parameters. To tackle this issue, we analyze OFT through the lens of information transfer and distill several key principles that underpin improved parameter efficiency. Drawing inspiration from the Cooley-Tukey fast Fourier transform, known for its optimized information propagation, we introduce a compact orthogonal parameterization based on butterfly structures. Integrating this mechanism into OFT yields Orthogonal Butterfly~(BOFT), a novel finetuning method that balances generality and parameter efficiency. Notably, BOFT generalizes OFT, incorporating it as a specific instantiation within its broader framework. We conduct comprehensive empirical analyses, adapting state-of-the-art vision transformers, large language models, and text-to-image diffusion models to an array of tasks in vision and language domains.
\end{abstract}

\section{Introduction}

Contemporary models, including ChatGPT~\cite{brown2020language,chen2021evaluating} and Stable Diffusion~\cite{rombach2022high}, illustrate the powerful generalization achieved by modern foundation models. This surge in performance correlates with substantial growth in model size

(\emph{e.g.}, GPT-3 comprises nearly 175B parameters), resulting in heightened barriers to training models from the ground up. Consequently, significant advancements hinge on efficient strategies for repurposing these models to new objectives, circumventing the need for exhaustive retraining. Three primary adaptation strategies exist: \emph{model finetuning}~\cite{hulora2022,zaken2022bitfit,guo2020parameter,raffel2020exploring,qiu2023controlling,zhang2022adaptive,chavan2023one}, which involves adjusting only a portion of the model parameters while maintaining the original architecture; \emph{adapter tuning}~\cite{rebuffi2017learning,houlsby2019parameter,pfeiffer2020adapterfusion,lin2020exploring,he2021towards}, where additional, dedicated trainable modules are incorporated; and \emph{prompt tuning}~\cite{li2021prefix,lester2021power}, which augments inputs with optimizable prefix tokens. Among these, model finetuning stands out for its simplicity, strong practical performance, and lack of impact on inference speed.

The central concept underlying model finetuning is to maintain a degree of similarity between the pretrained and finetuned models according to specific metrics, thereby retaining the knowledge acquired during pretraining. Traditional finetuning techniques, for example, frequently utilize a reduced learning rate; this heuristic strategy serves to limit the divergence between the pretrained model and its finetuned counterpart. Considering that weight matrices are inherently structured, a more systematic methodology aims to maintain the relational structure of the weights—specifically, the angles between neuron pairs. Building on this idea, the Orthogonal Finetuning~(OFT) framework~\cite{qiu2023controlling} has been proposed. In OFT, an orthogonal matrix is applied uniformly to neurons within the same layer, guaranteeing the preservation of pairwise angular relationships between neurons throughout the finetuning stage. While OFT has exhibited strong generalization and robust convergence properties, particularly when applied to finetuning text-to-image diffusion models~\cite{qiu2023controlling}, it faces challenges related to the sheer number of trainable parameters resulting from the large size of the orthogonal matrices involved. To mitigate this, OFT employs a block-diagonal structure, which significantly reduces parameter count. Nevertheless, this parameter reduction imposes limitations: the sparsity pattern of the orthogonal matrix becomes fixed, and each block undergoes orthogonal transformation independently. Such block-diagonal sparsity, despite its parameter efficiency, may inadvertently introduce restrictive inductive biases (for example, diminishing the matrix's expressiveness and hindering its capacity to replicate conventional linear transformations), and, crucially, optimal strategies for selecting the sparsity pattern are still unknown.

The fundamental challenge addressed here is constructing a dense orthogonal matrix in a parameter-efficient manner. At first, this appears difficult, since a fully dense orthogonal matrix of dimension $d$ conventionally needs $\mathcal{O}(d^2)$ parameters. However, our approach innovatively assembles a dense orthogonal matrix through the composition of multiple factorized sparse matrices. The key idea is to balance parameter count and computational effort: by opting for a representation where the orthogonal matrix is the product of sparse matrices, we reduce the number of trainable parameters at the cost of performing repeated matrix multiplications during each training cycle. We frame this matrix factorization as an information transmission problem—creating a dense orthogonal matrix via sparse matrix products becomes analogous to transmitting information through a grid-structured graph. This perspective reveals a flexible set of matrix factorization schemes that manage trainable parameter complexity while maintaining the ability to represent dense matrices.

To further maximize parameter efficiency within this information transmission view, we look to butterfly structures exemplified by the Cooley-Tukey fast Fourier transform algorithm~\cite{cooley1965algorithm}, where information can be propagated between nodes in a highly efficient manner~\cite{klasing1994broadcasting}. The butterfly graph intrinsic to the Cooley-Tukey method naturally suggests a form of sparse matrix decomposition, commonly known as butterfly factorization. Using this construction, we generate a $d$-dimensional dense matrix by sequentially multiplying $\mathcal{O}(\log d)$ sparse matrices, each featuring only $\mathcal{O}(d)$ nonzero entries. By maintaining orthogonality in each factor, we achieve an efficient orthogonal parametrization that requires just $\mathcal{O}(d\log d)$ parameters, dramatically less than conventional dense parameterizations. Building on this structure, we introduce Orthogonal Butterfly (BOFT), a new parameter-efficient finetuning strategy. BOFT generalizes OFT by providing a broad orthogonal finetuning framework. Notably, both the butterfly and block-diagonal designs are rooted in sparsity, which helps decrease the number of trainable parameters in orthogonal matrices. It is worth noting the relationship to LoRA~\cite{hulora2022}, which leverages low-rank structure—both low-rank and sparse matrices are prominent classes of structured matrices~\cite{candes2011robust} widely used to promote parameter efficiency.

In contrast to the block-diagonal approach adopted by OFT, which seeks a balance between expressivity and regularity, BOFT leverages the butterfly architecture to enable a more continuous spectrum between matrices from the full orthogonal group (which supports expressivity) and identity matrices (which support regularity). This methodology allows for the discovery of a more constrained hypothesis class within the orthogonal group tailored to downstream applications. Given that the butterfly structure underpins many efficient linear transforms—such as the discrete Fourier transform and the discrete cosine transform—our structured method is expected to introduce a form of implicit inductive bias that enhances generalization while mitigating the risk of overfitting. This rationale is anchored in the fact that butterfly structures can inherently represent a variety of established linear transforms, whereas the block-diagonal structures in OFT are unable to represent any of them. Our key contributions can be summarized as follows:

\begin{itemize}[leftmargin=*,nosep]
    \item We reframe the issue of parameter-efficient orthogonal finetuning using a novel information transmission perspective, which recasts the problem of constructing a dense, parameter-efficient orthogonal matrix as an information flow challenge over a grid-structured graph.
    \item Drawing from the butterfly structures in the Cooley-Tukey algorithm, we introduce Orthogonal Butterfly, a new parameter-efficient orthogonal finetuning scheme, set within the information transmission framework.
    \item We offer several theoretical justifications for BOFT's capacity to markedly reduce the number of learnable parameters without compromising expressivity or stability during training. The matrix factorization in BOFT also enables an intriguing interpolation between weights (refer to Figure~\ref{fig:boft_interpolation}).
    \item For the first time, we demonstrate the application of orthogonal finetuning beyond the context of controllable text-to-image synthesis~\cite{qiu2023controlling}, thereby highlighting its viability as a general-purpose model finetuning strategy. Specifically, we assess BOFT on a suite of adaptation scenarios spanning computer vision and natural language processing, showing that it achieves superior performance compared to leading methods, underscoring its remarkable parameter efficiency and capacity for generalization.
\end{itemize}

\section{Related Work}

\textbf{Parameter-efficient finetuning (PEFT)}. The continual growth in the size and capability of foundation models (\emph{e.g.}, \cite{brown2020language,radford2021learning,rombach2022high,kirillov2023segment}) has intensified interest in developing parameter-efficient approaches for adapting these models to new tasks~\cite{treviso2023efficient,ding2023parameter,lialin2023scaling}. Within the diverse range of PEFT techniques~\cite{houlsby2019parameter,aghajanyan2020intrinsic,hulora2022,edalati2022krona,wang2022adamix,gheini2021cross,zaken2022bitfit,guo2020parameter,sung2021training,ansell2022composable,lester2021power,li2021prefix,vu2022spot,he2021towards,mao2021unipelt,karimi2021compacter,liu2022few,sung2022lst,chen2022parameter,jia2022visual,chen2022adaptformer,zhang2022neural,jie2023fact,lian2022scaling,luo2023towards}, those based on reparameterization~\cite{aghajanyan2020intrinsic,hulora2022,edalati2022krona} bear the closest relationship to our method. LoRA~\cite{hulora2022}, for example, finetunes the weights of pretrained models by injecting a low-rank update formulated as the product of two low-rank matrices, achieving strong results in natural language applications. While LoRA enforces a uniform rank for these matrices across all layers, more recent efforts~\cite{zhang2022adaptive,valipour2022dylora,zhang2023increlora} have proposed adapting the rank on a per-layer basis to optimize the distribution of parameters. The simplicity of low-rank weight updates has contributed to their wide adoption~\cite{dettmers2023qlora,zi2023delta,chavan2023one}. Building on ideas from hyperspherical energy and its relationship to model generalization~\cite{liu2018learning,liu2021orthogonal}, \cite{qiu2023controlling} introduced orthogonal finetuning (OFT), a method that finetunes text-to-image diffusion models by learning an orthogonal transformation within a layer. OFT not only achieves enhanced generalization but also exhibits greater training stability in comparison to LoRA, although this often comes at the cost of using more trainable parameters. Thus, further improving the parameter efficiency of OFT presents a valuable research direction. Additionally, it remains unexplored whether OFT can generalize beyond text-to-image diffusion model adaptation~\cite{qiu2023controlling} to more diverse downstream tasks. The proposed BOFT addresses OFT's parameter cost by employing butterfly factorization, which allows us to reveal the strengths of orthogonal finetuning across a broader spectrum of adaptation scenarios.

\textbf{Butterfly structures}. The radix-2 Cooley-Tukey algorithm efficiently decomposes the $N$-point discrete Fourier transform into two $\frac{N}{2}$-point discrete Fourier transforms in a recursive manner, yielding a butterfly pattern that is representable as a product of sparse matrices—collectively termed a butterfly matrix. Applications of butterfly matrices include orthogonal matrix parameterization for streamlining Gaussian elimination~\cite{parker1995random}, stabilization of recurrent neural network training~\cite{jing2017tunable}, and kernel approximation~\cite{munkhoeva2018quadrature}. Work such as~\cite{dao2019learning,dao2019kaleidoscope} has demonstrated how to learn efficient linear transforms using butterfly parameterizations, while~\cite{chen2022pixelated,dao2022monarch} leverage them to accelerate neural network training. Furthermore, butterfly structures support fast matrix-vector computations~\cite{michielssen1996multilevel,o2010algorithm}, efficient data-sparse matrix representations~\cite{li2015butterfly}, and robust network communication~\cite{klasing1994broadcasting,li2003linear,rajasingh2016transmission}. Distinct from earlier studies, our work employs butterfly structures specifically to improve the parameter efficiency of OFT.

\section{An Information Transmission View on Orthogonal Finetuning}

We begin by reviewing fundamental aspects of OFT. In OFT, the finetuning process of a pretrained weight matrix involves reparameterizing it as the product of a trainable orthogonal matrix and the original frozen weight matrix. Unlike LoRA, which enhances weights via an additive low-rank update, OFT employs a multiplicative update using an orthogonal matrix. While LoRA's parameter efficiency is rooted in its low-rank design, classic OFT attains similar efficiency through a block-diagonal orthogonal configuration~\cite{qiu2023controlling}, where reducing block sizes leads to increased efficiency. Figure~\ref{fig:comp_lora} visually summarizes the comparison. The rationale for utilizing orthogonal transformations in finetuning is to maintain the angles between neurons~\cite{liu2018learning,liu2021orthogonal,qiu2023controlling}, thereby conserving much of the semantic information acquired during pretraining.

Specifically, OFT learns an orthogonal matrix $\bm{R}\in\mathbb{R}^{d\times d}$ alongside a pretrained linear layer $\bm{W}^0\in\mathbb{R}^{d\times n}$ and modifies the computation in the forward pass from $\bm{z}=(\bm{W}^0)^\top\bm{x}$ to $\bm{z}=(\bm{R}\bm{W}^0)^\top\bm{x}$, where $\bm{x}\in\mathbb{R}^{d}$ serves as input and $\bm{z}\in\mathbb{R}^{n}$ as output. Zero initialization in OFT is achieved by setting $\bm{R}$ to the identity matrix at the outset. To preserve orthogonality during training, we adopt the Cayley parameterization, consistent with prior work~\cite{qiu2023controlling,liu2021orthogonal}: $\bm{R}=(\bm{I}+\bm{Q})(\bm{I}-\bm{Q})^{-1}$, where $\bm{Q}$ denotes a skew-symmetric matrix satisfying $\bm{Q}=-\bm{Q}^\top$. 

For parameter efficiency, the orthogonal matrix $\bm{R}$ is structured in a block-diagonal fashion, parameterized as $\text{diag}(\bm{R}_1,\bm{R}_2,\cdots,\bm{R}_r)$ with each block $\bm{R}_i\in\mathbb{R}^{b\times b}$, such that $br=d$. However, this design presupposes that a neuron's dimensions (i.e., the elements of each column in $\bm{W}^0$) can be arbitrarily split into $r$ separate groups, with each group transformed independently via distinct orthogonal blocks. Although such block-diagonal schemes have shown empirical benefit, they rest on the somewhat artificial assumption of grouping dimensions by index, which limits flexibility. Therefore, densely-parameterized orthogonal matrices would be advantageous. This consideration prompts the question: \emph{Is it possible to construct a dense orthogonal matrix while preserving parameter-efficiency?}

To tackle this problem, we suggest constructing a dense orthogonal matrix as the product of several sparse orthogonal matrices. In order to systematically and intuitively analyze the sparsity structures in orthogonal matrix decomposition, we conceptualize the process of generating a dense orthogonal matrix in OFT as an information transmission scenario. More concretely, forming a dense matrix $\bm{R}\in\mathbb{R}^{d\times d}$ through multiplication of $m$ square matrices, $\bm{R}=\bm{B}_m\bm{B}_{m-1}\cdots\bm{B}_1$, can be interpreted as relaying information through a grid comprising $d\times (m+1)$ vertices, as depicted in Figure~\ref{fig:info_trans}. This analogy arises from the perspective that a dense square matrix of dimension $d$ characterizes comprehensive connectivity between two sets of $d$ nodes. Here, if an entry $R_{ij}$ in $\bm{R}$ equals zero, it signifies that communication from node $j$ to node $i$ is not possible; conversely, a nonzero $R_{ij}$ permits information flow. Accordingly, representing the dense matrix $\bm{R}$ as a product of matrices $\bm{B}_m\bm{B}_{m-1}\cdots\bm{B}_1$ can be viewed as a sequential exchange of information governed by the graphs generated by each $\bm{B}_i$. The transfer of information proceeds according to the connectivity prescribed first by $\bm{B}_1$ and concludes with $\bm{B}_n$. As a concrete illustration in Figure~\ref{fig:info_trans}, we study the decomposition $\bm{R}=\bm{B}_5\bm{B}_4\bm{B}_3\bm{B}_2\bm{B}_1$ and visualize both the associated sparsity patterns and the corresponding graph. The depicted graph results from unfolding the matrix product step by step. Within this constructed graph, each $\bm{B}_i$ encodes the connectivity from nodes in layer $i$ to those in layer $i+1$. Precisely, an entry $(j_1,j_2)$ in $\bm{B}_i$ indicates whether there exists a directed connection from the $j_2$-th node in the $i$-th layer to the $j_1$-th node in the subsequent $(i+1)$-th layer (a zero indicating absence of connection). In order for $\bm{B}_5\bm{B}_4\bm{B}_3\bm{B}_2\bm{B}_1$ to be fully dense, information must propagate from every node in the first layer to every node in the sixth layer. Should we restrict attention to $\bm{R}=\bm{B}_4\bm{B}_3\bm{B}_2\bm{B}_1$—that is, connections between the initial layer and the fifth layer—we might observe, for example, that node $1$ is unable to transmit to node $3$. Consequently, the sparsity structure of $\bm{B}_4\bm{B}_3\bm{B}_2\bm{B}_1$ reflects a zero in position $(3,1)$. Similarly, for factorizations $\bm{R}=\bm{B}_3\bm{B}_2\bm{B}_1$ or $\bm{R}=\bm{B}_2\bm{B}_1$, the induced graph faithfully corresponds to the underlying sparsity patterns. In general, for a matrix $\bm{R}\in\mathbb{R}^{d\times d}$ to attain density, the collection of $m$ factor matrices $\bm{B}_m,\cdots,\bm{B}_1$ should define a series of directed connections over a $d\times (m+1)$ grid, in which each directed edge links two nodes residing in consecutive layers (or columns), ensuring comprehensive information transfer from every source node in the first layer to every destination node in the final $(m+1)$-th layer.

The information transmission perspective provides valuable insights into the sparsity characteristics of factorization matrices within OFT. In Figure~\ref{fig:blk_diag}, the block-diagonal configuration of $\bm{R}$ from standard OFT is depicted. While this design decreases the total number of learnable parameters, it does not permit $\bm{R}$ to be a dense matrix. The central objective is to construct a dense orthogonal matrix using $m$ sparse orthogonal matrices, optimizing for a minimal count of effective learnable parameters. When viewed through the lens of information transmission, two main principles emerge: \emph{(i)} \textbf{dense connectivity}, meaning that each node in the initial layer is connected to every node in the final layer via at least one path, and \emph{(ii)} \textbf{minimal free edges}, meaning the total number of edges remains as low as possible without violating orthogonality. The orthogonality condition introduces nuanced restrictions on connections between adjacent layers. For instance, for each matrix $\bm{B}_i$ to be of full rank—a prerequisite for orthogonality—it is necessary to establish $d$ edges to enable a bijective correspondence between all nodes in level $i$ and level $(i+1)$. As a result, the minimum number of edges required between any two successive layers is $d$ (such as 4, as illustrated in Figure~\ref{fig:info_trans}). These essential $d$ edges are dictated by orthogonality and are not part of the trainable parameter count, since their values are fixed (for a $d\times d$ orthogonal matrix with $d$ non-zeros, these elements are restricted to $\pm 1$). Owing to the fact that orthogonal matrices need fewer parameters than general matrices, enforcing orthogonality leads to extra dependencies among the edges. For example, in a $2\times 2$ orthogonal matrix, a zero at the $(1,1)$ entry necessitates a zero at the $(2,2)$ entry as well (that is, omitting one edge can force the omission of another). Thus, orthogonality may result in automatic addition or removal of certain edges, contingent on the allowed connections. The information transmission view is especially instructive in the context of OFT when examining the nonzero structure of the resulting orthogonal matrix, as the precise values of the trainable nonzero entries in $\bm{R}$ are inconsequential—only their locations matter.

A straightforward fully-connected configuration between two layers necessitates $\mathcal{O}(d^2)$ connections (\emph{i.e.}, represented by one dense orthogonal matrix), resulting in $d^2-d$ individual edges (for instance, with $d=4$, this amounts to 12 connections). As depicted in Figure~\ref{fig:info_trans}, one can construct a viable matrix factorization that uses a total of $10$ edges, which is in fact fewer than required by a single dense orthogonal matrix. This conceptual framework allows for an examination of OFT's parameter-efficiency through a graphical lens, making it straightforward to devise suitable factorizations. Our methodology is influenced by the distinctive topology of butterfly graphs from the Cooley-Tukey algorithm~\cite{cooley1965algorithm}, which are capable of densely connecting $d$ source and $d$ target nodes using only $\mathcal{O}(d\log d)$ edges. For instance, the configuration in Figure~\ref{fig:info_trans} achieves full connectivity with $10$ connections, while the butterfly network topology achieves the same with just $8$ edges. In the subsequent section, we demonstrate how leveraging the butterfly structure can lead to improved parameter efficiency.

\section{Orthogonal Parameterization by Butterfly Factorization}

The butterfly structure originates from its application in the Cooley-Tukey algorithm for efficient computation of the fast Fourier transform. Within the context of Fourier transforms, a modification in the frequency domain induces widespread changes throughout the spatial domain—an analogy to our problem of information dissemination, where every node in the initial layer should be capable of relaying information to all terminal-layer nodes. Additionally, the butterfly pattern is a well-established topology for communication networks~\cite{solihin2015fundamentals,li2011linear} designed for streamlined data transfer. For $k\geq 2$, where $k$ is a power of $2$, we introduce the \emph{butterfly factor} $\bm{B}^F(k)$ as follows:

\begin{equation}\label{eq:but_factor}
\begin{aligned}
    \bm{B}^F(k)=\begin{bmatrix}
    \text{diag}(\bm{d}_1) & \text{diag}(\bm{d}_2) \\
    \text{diag}(\bm{d}_3) & \text{diag}(\bm{d}_4)
    \end{bmatrix}\in\mathbb{R}^{k\times k},
\end{aligned}
\end{equation}

where $\bm{d}_i\in\mathbb{R}^{\frac{k}{2}}$ for $i=1,2,3,4$ represent specific vectors. When $d=2^N$, we further define the $d$-dimensional \emph{butterfly matrix} $\bm{B}(d)\in\mathbb{R}^{d\times d}$ recursively as:

\begin{equation}
\begin{aligned}
    \!\!\bm{B}(d)=\tilde{\bm{B}}(d,d)\!\cdot\!\begin{bmatrix}
    \bm{B}_1(\frac{d}{2})\!\!\!\!\! & \bm{0} \\
    \bm{0}\!\!\!\!\! &\bm{B}_2(\frac{d}{2})
    \end{bmatrix}= \tilde{\bm{B}}(d,d)\tilde{\bm{B}}(d,\frac{d}{2})\cdots\tilde{\bm{B}}(d,2),
\end{aligned}
\end{equation}

Here, $\bm{B}_1(\frac{d}{2})$ and $\bm{B}_2(\frac{d}{2})$ denote butterfly matrices of dimension $\frac{d}{2}$. We introduce the notion of a \emph{butterfly component}, denoted as $\thickmuskip=2mu \medmuskip=2mu \tilde{\bm{B}}(d,k)=\text{diag}(\bm{B}^F_1(k),\cdots,\bm{B}^F_{d/k}(k))$, which is a block-diagonal matrix with size $d \times d$, where each block is of size $k$. In this formulation, each $\bm{B}^F_i(k)$, for all $i$, represents a butterfly factor as detailed in Equation~\ref{eq:but_factor}. We now outline the use of butterfly matrices as a means to parameterize orthogonal matrices. This construction requires that all elements $\tilde{\bm{B}}(d,k)$ comprising the product in the butterfly matrix $\bm{B}(d)$ are themselves orthogonal for all $k$. 

To illustrate, consider $\tilde{\bm{B}}(d,2)$, a block-diagonal matrix whose blocks are of size $2$. By parameterizing each $2\times 2$ block using the Cayley transform (or equivalently, a $2$-dimensional rotation), as described in \cite{liu2021orthogonal,qiu2023controlling}, $\tilde{\bm{B}}(d,2)$ is guaranteed to be orthogonal. Each butterfly component’s sparsity structure results from a permutation applied to the sparsity pattern in $\tilde{\bm{B}}(d,2)$, enabling us to represent all butterfly components as orthogonal matrices through a straightforward extension. This yields an efficient orthogonal parameterization composed of multiple $2\times2$ orthogonal matrices. 

Building on \cite{chen2022pixelated}, we extend the butterfly matrices to define a block butterfly component $\tilde{\bm{B}}^b(d,k)$, where every element in $\bm{d}_i$ for all $i$ is replaced by a $b\times b$ matrix. For the case $k=2$, ensuring orthogonality of the block butterfly component $\tilde{\bm{B}}^b(d,2)$ is accomplished by parameterizing each $2b\times 2b$ block as an orthogonal matrix. The nonzero structure of any other block butterfly component $\tilde{\bm{B}}^b(d,k)$ with $k>2$ is given by block-wise permutations of the nonzero pattern in $\tilde{\bm{B}}^b(d,2)$, and thus they too can be orthogonalized using the same approach. As a result, the forward computation in BOFT proceeds as follows:

\begin{equation}
    \begin{aligned}\nonumber
    \bm{z} = \big(\bm{R}(m, b)\cdot \bm{W}^0\big)^\top \bm{x}, ~~~~ \text{s.t.}~ \bigg\{\bm{R}(m, b) = \prod_{i=1}^m \tilde{\bm{B}}^b_{(d,i)} ~~\&~~ \underbrace{\big(\tilde{\bm{B}}^b_{(d,j)}\big)^\top \tilde{\bm{B}}^b_{(d,j)} = \tilde{\bm{B}}^b_{(d,j)} \big(\tilde{\bm{B}}^b_{(d,j)}\big)^\top \!= \bm{I}_d}_{\forall j \in [1, m]}\bigg\},
    \end{aligned}
\end{equation}

Here, for notational clarity, $\tilde{\bm{B}}^b(d, 2^{m-i+1})$ is simply written as $\tilde{\bm{B}}^b_{(d,i)}$, and $\bm{I}_d$ stands for the $d$-dimensional identity matrix. The orthogonal matrix $\bm{R}(m, b) \in \mathbb{R}^{d \times d}$ is constructed as a product of multiple orthogonal butterfly blocks. We abbreviate BOFT with $\bm{R}(m, \frac{b}{2})$ as BOFT($m$, $b$) when $b \geq 2$. Notably, if $\thickmuskip=2mu \medmuskip=2mu m=1$, then BOFT($1$, $b$) corresponds to the block-diagonal OFT~\cite{qiu2023controlling} of block size $b$. In the special case where $b=d$, BOFT($1$, $d$) collapses to the original OFT~\cite{qiu2023controlling} parameterized by an unrestricted orthogonal matrix. For BOFT($\log \frac{2d}{b}$, $b$), the construction leverages the block butterfly matrix $\bm{B}^{\frac{b}{2}}(d)$ as $\bm{R}$, resulting in a dense orthogonal matrix. Broadly, BOFT($m$, $b$) introduces a total of $\frac{1}{2}(b-1)dm$ tunable parameters when fine-tuning a linear mapping of size $d \times n$. Employing a standard butterfly matrix, i.e., setting $m = \log d, b = 2$, restricts BOFT to $\mathcal{O}(d \log d)$ parameters. In comparison, the vanilla OFT configured with a dense orthogonal matrix demands $\mathcal{O}(d^2)$ parameters, whereas a block-diagonal OFT with $r$ blocks utilizes $\mathcal{O}(bd)$. Consequently, to achieve a dense orthogonal transformation, the original OFT requires a block size of $b = d$, whereas BOFT enables such transformations with arbitrary $b$ values.

\textbf{Identity Initialization in BOFT}. In finetuning scenarios, it is common to initialize with the exact parameters from the pretrained model, to ensure that the subsequent updates introduce only minor deviations. For instance, LoRA initializes the low-rank component with zeros. In the case of BOFT, each butterfly block is initialized as an identity matrix (that is, the corresponding skew-symmetric matrices in the Cayley transform are set to zero).

\textbf{Multiplicative Dropout for BOFT}. LoRA~\cite{hulora2022} augments the low-rank adaptation by introducing a Dropout layer to curb overfitting. Whereas traditional Dropout~\cite{srivastava2014dropout} is compatible with LoRA, the multiplicative nature of the weight modification in BOFT necessitates a modified approach. To this end, we propose a multiplicative Dropout variant tailored for BOFT. Since $\bm{R}(m, b)$ is an aggregation of $m$ butterfly units, which can themselves be rearranged into $2b \times 2b$ block-diagonal orthogonal forms, the procedure first randomly selects $p_1$ percent of butterfly units and, within each, $p_2$ percent of diagonal blocks. These selected blocks are then replaced with identity matrices during Dropout.

\section{Intriguing Insights and Discussions}

\textbf{Expressivity of BOFT}. The butterfly architecture, when combined with appropriate permutations, is capable of exactly representing a variety of standard fast linear transforms~\cite{dao2019learning,dao2019kaleidoscope} (\emph{e.g.}, the fast Fourier transform, Hadamard transform). Nevertheless, the degree to which our orthogonal butterfly matrix can serve as an accurate surrogate for an arbitrary orthogonal matrix has not been thoroughly investigated. To shed light on this, we carry out an experiment aiming to approximate a random dense orthogonal matrix~\cite{anderson1987generation} of size $512\times 512$. The results in Figure~\ref{fig:para_eff} are computed as averages over 10 independent trials. The vertical axis represents the approximation error, while the horizontal axis indicates the count of effective, trainable parameters. Within each colored curve, BOFT is evaluated using a constant block size; the leftmost data point corresponds to BOFT($1$,$b$), i.e., the conventional block-diagonal OFT using block size $b$. Generally, BOFT offers greater parameter efficiency compared to OFT. For instance, BOFT($9$,$2$) exhibits higher expressive power than BOFT($1$,$16$), yet requires significantly fewer parameters. Using smaller values of $b$ in conjunction with larger $m$ further enhances parameter efficiency; as an example, BOFT($6$,$4$) achieves comparable approximation error to BOFT($2$,$16$), but with a substantial reduction in parameter count. Overall, the butterfly structure realizes a more constrained and organized subset of the orthogonal group than block-diagonal matrices, providing theoretical guarantees that BOFT surpasses OFT in expressivity at a fixed block size.

\begin{theorem}[Expressivity of BOFT]\label{thm:exp_boft}
    BOFT exhibits superior expressivity compared to OFT when using the same block size. To ensure that a butterfly matrix is capable of expressing any orthogonal matrix of dimension $d$, it suffices to form products of the type $\bm{B}_{d-1,1}(d)\bm{B}^\top_{d-1,2}(d)\cdots\bm{B}_{1,1}(d)\bm{B}^\top_{1,2}(d)$, with all $\bm{B}_{i,j}(d)$ denoting butterfly matrices.
\end{theorem}

The assertion in Theorem~\ref{thm:exp_boft} motivates a natural extension for BOFT: the overall orthogonal matrix may be expressed as $\thickmuskip=2mu \medmuskip=2mu \bm{R}^G(m_1,b_1,m_2,b_2,l)=\bm{R}_{l,1}(m_1,b_1)\bm{R}_{l,2}^T(m_2,b_2)\cdots\bm{R}_{1,1}(m_1,b_1)\bm{R}_{1,2}^T(m_2,b_2)$, where $\bm{R}_{i,j}^T(m,b)$ denotes the orthogonal matrices deployed in BOFT. Specifically, for $m_1=m_2=\log d$, $b_1=b_2=2$, and $l=d-1$, $\bm{R}^G(m_1,b_1,m_2,b_2,l)$ is sufficiently expressive to parameterize the entirety of the orthogonal group. This compositional pattern is also referred to as the kaleidoscope hierarchy~\cite{dao2019kaleidoscope}. It is important to emphasize, however, that greater expressive power does not necessarily translate into improved performance during finetuning. In practice, although full finetuning theoretically enables arbitrary expressivity, it often leads to suboptimal outcomes. Balancing the expressivity of the model with its structural regularization remains central for achieving strong generalization during finetuning. By introducing structured priors, BOFT expands the feasible parameter space, which facilitates the search for a more optimal trade-off between flexibility and regularity.

\textbf{Spectral properties}. In comparison to LoRA, orthogonal finetuning typically achieves superior spectral characteristics, as it exactly maintains the spectral norm of the pretrained parameter matrix $\bm{W}^0$. This can be observed through singular value decomposition: $\bm{W}^0=\bm{U}\bm{\Sigma}\bm{V}^\top$, with $\bm{U},\bm{V}$ as orthogonal matrices and $\bm{\Sigma}$ as the diagonal matrix containing singular values. With OFT and BOFT, the procedure involves multiplying by an orthogonal matrix $\bm{R}$ from the left, resulting in finetuned weights of the form $\bm{R}\bm{U}\bm{\Sigma}\bm{V}^\top$. This operation leaves the maximal singular value—equivalently, the spectral norm of $\bm{W}^0$—unaltered. Such norm preservation has been demonstrated to significantly promote both stability during training and model generalization~\cite{miyato2018spectral,yoshida2017spectral}. Additional mathematical discussions can be found in Appendix~\ref{app:sn_property}.

\textbf{Orthogonal finetuning as learning bilinear similarity}. BOFT is interpretable as a process of learning the bilinear similarity $\bm{w}^0_i\bm{R}\bm{x}$, where $\bm{w}^0_i$ denotes the $i$-th column (i.e., neuron) of the weight matrix $\bm{W}^0$. Within this framework, BOFT aims to learn the similarity matrix $\bm{R}$, subject to a strict constraint of orthogonality. This approach is inherently related to distance metric learning~\cite{xing2002distance} as well as bilinear forms~\cite{roman2005advanced}.

\textbf{Inductive bias and generalization in BOFT}. The matrix $\bm{R}(m,b)$ employed in BOFT is typically structured, representing a specific subset of the orthogonal group and thus restricting the hypothesis space. Consequently, BOFT imposes an inductive bias. We posit that the inductive bias introduced by the butterfly factorization is advantageous for model generalization, as its shared structural pattern aligns with numerous well-known linear transforms~\cite{dao2019kaleidoscope} including the discrete Fourier, sine/cosine, and Hadamard transforms. Additionally, the sparse factorization inherent to BOFT may introduce further implicit inductive biases~\cite{gunasekar2017implicit,li2018algorithmic,liu2019neural}.

\textbf{Contrast with butterfly-based sparse training}. Numerous prior studies~\cite{dao2019learning,dao2019kaleidoscope,chen2022pixelated,dao2022monarch} have explored sparse training schemes involving the butterfly parameterization. Typically, these works focus on reparameterizing neural network weight matrices directly with butterfly structures and initiate training from scratch. In a departure from this, \cite{dao2022monarch} explores the finetuning of pretrained models by projecting their weights onto modified butterfly matrices, subsequently updating the projected parameters for transfer to downstream applications. BOFT introduces a markedly distinct approach to finetuning by modifying weights through weight matrices that are shared across layers.

\input{rebuttal-results}

\section{Concluding Remarks and Limitations}

This paper presents Orthogonal Butterfly, an efficient and broadly applicable finetuning technique for foundation models grounded in the butterfly matrix framework. Our main contribution toward greater parameter efficiency lies in expressing a dense orthogonal transformation as a product of several sparse orthogonal matrices. To facilitate the discovery of appropriate matrix decompositions, we introduce a graphical framework for modeling information propagation. Within this setting, we identify the butterfly configuration as particularly effective for achieving the goals of sparse orthogonal factorization. Empirical results demonstrate the practical advantages of BOFT when adapting large language models, extensive vision architectures, and text-to-image generation systems. Furthermore, our evaluation provides evidence for BOFT’s effectiveness as a general-purpose finetuning paradigm.

Nonetheless, there are inherent limitations to BOFT. Because the ultimate orthogonal matrix in BOFT is realized as a product of several orthogonal matrices, training involves slightly greater computational overhead relative to OFT. Addressing the runtime efficiency of BOFT during training remains a challenge for future work. Notably, however, after finetuning is complete, the orthogonal matrices produced by BOFT can be folded into the existing pretrained weights, resulting in no extra inference-time latency. Additionally, it remains unclear if the butterfly architecture constitutes the most optimal approach for information transfer. Our proposed framework for information flow further suggests that lessons can be drawn from the domain of computer networking, where various topological structures are evaluated for communication efficiency. We anticipate that alternative network designs may further enhance the construction of dense orthogonal matrices.